%% =====================================================================
%%  Supplementary Material for:
%%  Confidence-Gated Dual-Process Fall Detection (RCDP), IEEE JBHI
%%
%%  Build separately:  tectonic supplementary.tex  (or pdflatex x2)
%%  Upload the PDF as "Supplementary Material" in ScholarOne.
%%  Table order is fixed (S1--S9); main.tex cites these numbers literally.
%%  Do not reorder the \input lines without updating main.tex.
%% =====================================================================
\documentclass[journal]{IEEEtran}

\usepackage{amsmath,amssymb,amsfonts}
\usepackage{booktabs}
\usepackage{array}
\usepackage{multirow}
\usepackage{url}
\usepackage{xcolor}
\usepackage[hidelinks]{hyperref}

\newcommand{\authorcheck}[1]{}

\renewcommand{\thetable}{S\arabic{table}}
\renewcommand{\thesection}{S-\Roman{section}}
\renewcommand{\thefigure}{S\arabic{figure}}

\begin{document}

\title{Supplementary Material for\\ ``Confidence-Gated Dual-Process Fall Detection: Retrieval-Augmented Reasoning and Cost-Sensitive Response on Wearable Inertial Sensor Data''}
\author{\authorcheck{Same author list as the main paper.}}
\maketitle

Table and section numbers prefixed with ``S'' refer to this document; all other numbers refer to the main paper.

\section{Reproducibility Notes}
\label{sup:repro}
The baseline mode is called \texttt{primary\_only} in the paper; released configurations and result files use the legacy key \texttt{tier1\_only} (\texttt{configs/tier1\_sisfall.yaml}). Main classification results use the supplied main-ablation confusion counts, one observation per mode and fold; latency and escalation come from a separate timing run with identical TP/FN counts and FP differences of at most $2$ windows per fold. The shipped \texttt{configs/paper\_protocol.yaml} sets \texttt{primary\_backbone: cnn\_lstm\_attn} to match the main CNN--LSTM--Attention tables; HARMamba is a generality backbone only. The following scripts reproduce the tabulated arithmetic without a LaTeX installation:
\begin{verbatim}
python check_manuscript.py   # inputs, labels, cites, arity
python check_numbers.py      # Mean rows vs. per-fold rows
python check_reported_results.py # counts vs. tables
python recompute_ci.py       # paired fold intervals
\end{verbatim}

\section{Protocol Details}
\label{sup:protocol}
\paragraph{Budgets} Seed $42$; \texttt{max\_train}${}={}$40000, \texttt{max\_val}${}={}$5000, \texttt{max\_cal}${}={}$2000; main \texttt{max\_test}${}={}$1000 with a D18/D19 quota of $80$ windows/code; LLM comparison $150$ escalated cases/fold; feature adjudicators $160$/fold; costs $\mathrm{FN}{=}10$, $\mathrm{FP}{=}1$.
\paragraph{Gate grid} $\tau_{\mathrm{low}}$: $20$ points in $[0.05,0.45]$; $\tau_{\mathrm{high}}$: $20$ points in $[0.55,0.95]$; objective $\mathrm{FP}+\beta\,\mathrm{FN}+0.5\cdot(\text{deferred windows})$ over edge decisions; at least $20\%$ edge decisions (or $10$ windows) per candidate pair.
\paragraph{Reasoner} Ollama at \texttt{127.0.0.1:11434}; \texttt{mistral:latest} (primary), \texttt{qwen2.5-coder:7b-instruct} (comparison); temperature $0.1$; timeout $60$\,s; two retries; heuristic fallback. Output JSON keys: \texttt{prediction}$\in\{\texttt{fall},\texttt{adl}\}$, \texttt{severity}$\in\{\texttt{none},\texttt{mild},\texttt{moderate},\texttt{severe}\}$, \texttt{confidence}$\in[0,1]$, \texttt{rationale}, \texttt{action}. Prompts are in \texttt{prompts/}.
\paragraph{Action map (simulated)} Fall: impact $\ge 5\,g$ or stillness $\ge 10$\,s $\rightarrow$ emergency; impact $\ge 3\,g$ or stillness $\ge 3$\,s $\rightarrow$ notify; otherwise monitor. ADL: impact $\ge 2\,g$ or free-fall $\ge 0.12$\,s $\rightarrow$ monitor; otherwise log.

\section{Supplementary Results}
\label{sup:results}
Tables~S1--S9 report detector-only context (backbone zoo, seeds, sensor perturbation), gate operating points, per-fold KFall results, cost-ratio sensitivity, rationale grounding, escalation and latency. The robustness and seed runs use different test slices from the main protocol, so their absolute values differ slightly from the per-fold table of the main paper (Table~II).

\begin{table}[!t]
\centering
\caption{Five-fold detector-only SisFall context at the reported Youden operating points, $20{,}000$ standard-protocol windows/fold (prevalence $0.329$--$0.345$). E1 and retrained zoo are distinct checkpoint runs. Mean$\pm$sample std are recomputed from the supplied three-decimal fold F1 values; last-digit differences from separately supplied aggregate summaries can result from rounding. Neither run is the frozen main comparison or the separate $30{,}000$-window v2 evaluation. Scores describe tested implementations and training settings, not a universal architecture ranking.}
\label{tab:backbone}
\begin{tabular}{lcc}
\toprule
Model & E1 F1 & Retrained F1 \\
\midrule
CNN1D & 0.701$\pm$0.031 & 0.734$\pm$0.026 \\
LSTM & 0.693$\pm$0.016 & 0.711$\pm$0.026 \\
CNN--LSTM & 0.696$\pm$0.010 & 0.724$\pm$0.036 \\
CNN--LSTM--Attention & 0.717$\pm$0.022 & 0.739$\pm$0.030 \\
TCN & 0.702$\pm$0.029 & 0.747$\pm$0.034 \\
ModernTCN & 0.743$\pm$0.015 & 0.760$\pm$0.030 \\
TSMixer & 0.671$\pm$0.022 & 0.686$\pm$0.018 \\
PatchTST & 0.541$\pm$0.009 & 0.566$\pm$0.014 \\
Transformer & 0.690$\pm$0.013 & 0.703$\pm$0.017 \\
\bottomrule
\end{tabular}
\end{table}

\begin{table}[!t]
\centering
\caption{Multi-seed robustness on SisFall Fold~0 ($n{=}1000$, Ollama/Mistral backend). Seeds redraw the stratified test slice as well as training randomness, so absolute values differ from the Fold~0 row of the main-paper per-fold table (Table~II, which uses the locked seed-42 slice); the quantity of interest is the across-seed spread, not the level.}
\label{tab:seeds}
\setlength{\tabcolsep}{3pt}
\begin{tabular}{ccccc}
\toprule
Seed & Primary F1 & Stack F1 & Stack Rec. & Cost/1k \\
\midrule
42 & 0.747 & 0.909 & 0.998 & 95 \\
43 & 0.740 & 0.907 & 0.995 & 105 \\
44 & 0.746 & 0.899 & 1.000 & 96 \\
\midrule
Mean$\pm$std & 0.745$\pm$0.004 & \textbf{0.905$\pm$0.005} & 0.998$\pm$0.003 & 99$\pm$6 \\
\bottomrule
\end{tabular}
\end{table}

\begin{table}[!t]
\centering
\caption{Sensor-perturbation robustness of primary-detector routing under the cost gate (SisFall, five-fold mean F1). Rotations apply a random rigid rotation of the sensor frame with the stated bound, emulating garment or mounting drift; resampling emulates a different device sampling rate. Degradation is graceful---the worst case, $\pm15^\circ$, loses $0.014$ F1---so the ambiguous band is not an artefact of one exact sensor placement.}
\label{tab:robustness}
\begin{tabular}{lcc}
\toprule
Condition & Mean F1 & $\Delta$ vs.\ clean \\
\midrule
clean             & 0.768 $\pm$ 0.017 & --- \\
resample 100\,Hz  & 0.766 $\pm$ 0.016 & $-0.002$ \\
resample 50\,Hz   & 0.766 $\pm$ 0.015 & $-0.002$ \\
rotate $\pm10^\circ$ & 0.763 $\pm$ 0.020 & $-0.005$ \\
rotate $\pm15^\circ$ & 0.754 $\pm$ 0.023 & $-0.014$ \\
\bottomrule
\end{tabular}
\end{table}

\begin{table}[!t]
\centering
\caption{Archived pooled routing diagnostic without System~2. The entries labelled min-cost and Q25 tie at reported cost $1015$, while Q25 has $23$ percentage points less deferral. Under these displayed objectives, min-cost is dominated; the source selection rule and the treatment of deferred-case errors must be distinguished from the production routing surrogate. These pooled points do not justify choosing deployment thresholds without held-out calibration and adjudicator evaluation.}
\label{tab:pareto}
\begin{tabular}{lcccr}
\toprule
Operating point & $\tau_{\mathrm{low}}$ & $\tau_{\mathrm{high}}$ & Esc.\ rate & Cost \\
\midrule
Min cost            & 0.05 & 0.60 & 0.57 & 1015 \\
\textbf{Q25 esc.\ (dominant)} & \textbf{0.30} & \textbf{0.70} & \textbf{0.34} & \textbf{1015} \\
Q50 esc.            & 0.40 & 0.95 & 0.48 & 1042 \\
Q75 esc.            & 0.30 & 0.95 & 0.61 & 1223 \\
\bottomrule
\end{tabular}
\end{table}

\begin{table*}[!t]
\centering
\caption{Second-dataset evaluation on KFall, frozen primary detector versus full stack under the binary Fall/ADL protocol ($n{=}1000$/fold). This is not evidence of zero-shot transfer of SisFall-trained weights or formal non-inferiority. Fold~1 regresses; Fold~4 has the largest displayed F1 gain. Deltas below are computed from the displayed rounded values. Fold~4 primary row is consistent with its confusion counts (TP${=}473$, TN${=}508$, FP${=}9$, FN${=}10$; $483$ falls, not $\approx$499 as in other folds): F1 $0.980$, recall $0.979$, cost $109$.}
\label{tab:kfallperfold}
\begin{tabular}{c r ccc ccc cc}
\toprule
Fold & $n$ & \multicolumn{3}{c}{\texttt{primary\_only}} & \multicolumn{3}{c}{\texttt{gate\_knn\_llm}} & $\Delta$F1 & $\Delta$Cost \\
\cmidrule(lr){3-5}\cmidrule(lr){6-8}
 & & F1 & Rec. & Cost/1k & F1 & Rec. & Cost/1k & & \\
\midrule
0 & 1000 & 0.992 & 0.998 & 17  & 0.993 & 1.000 & 7  & $+0.001$ & $-10$ \\
1 & 1000 & 0.987 & 0.996 & 31  & 0.986 & 0.996 & 32 & $-0.001$ & $+1$  \\
2 & 1000 & 0.987 & 0.994 & 40  & 0.989 & 0.996 & 29 & $+0.002$ & $-11$ \\
3 & 1000 & 0.992 & 0.990 & 53  & 0.995 & 0.996 & 23 & $+0.003$ & $-30$ \\
4 & 1000 & 0.980 & 0.979 & 109 & 0.989 & 0.994 & 38 & $+0.009$ & $-71$ \\
\midrule
Mean & 1000 & 0.987 & 0.991 & 50 & \textbf{0.990} & \textbf{0.996} & \textbf{26} & $+0.003$ & $-24.2$ \\
\bottomrule
\end{tabular}
\end{table*}

\begin{table}[!t]
\centering
\caption{Sensitivity to the clinical cost ratio $\beta=c_{\mathrm{FN}}/c_{\mathrm{FP}}$ (SisFall, five-fold mean cost per $1000$ windows, $\mathcal{R}_\beta=\beta\cdot\mathrm{FN}+\mathrm{FP}$). The advantage of the stack is monotone in $\beta$: the more a deployment penalises missed falls relative to nuisance alarms, the larger the absolute saving. $\tau^\ast=1/(\beta{+}1)$ is the corresponding hard-decision Bayes threshold~\eqref{eq:tau_star}. The $\beta{=}10$ row is the paper's primary setting.}
\label{tab:betacost}
\begin{tabular}{ccl rr}
\toprule
$\beta$ & $\tau^\ast$ & Mode & Cost/1k & $\Delta$ \\
\midrule
5  & 0.167 & \texttt{primary\_only} & 628 & --- \\
5  & 0.167 & \textbf{\texttt{gate\_knn\_llm}} & \textbf{149} & $-479$ \\
\addlinespace
10 & 0.091 & \texttt{primary\_only} & 1151 & --- \\
10 & 0.091 & \textbf{\texttt{gate\_knn\_llm}} & \textbf{200} & $-951$ \\
\addlinespace
20 & 0.048 & \texttt{primary\_only} & 2197 & --- \\
20 & 0.048 & \textbf{\texttt{gate\_knn\_llm}} & \textbf{302} & $-1895$ \\
\bottomrule
\end{tabular}
\end{table}

\begin{table}[!t]
\centering
\caption{Rationale grounding. $\mathcal{G}_{\mathrm{faith}}$ is the fraction of rationales that mention at least one evidence field present in $\Sigma(\mathbf{x})$. It is a prompt-compliance check, not a measure of reasoning faithfulness or correctness: a rationale can cite a field and still reach the wrong conclusion. \emph{Graded cost} scores the four-level simulated action rather than the binary label. Contrastive critique with frozen labels raises $\mathcal{G}_{\mathrm{faith}}$ ($0.436{\rightarrow}0.843$) and lowers graded action cost by $45\%$.}
\label{tab:grounding}
\begin{tabular}{ccccc}
\toprule
\multirow{2}{*}{Fold} & \multicolumn{2}{c}{$\mathcal{G}_{\mathrm{faith}}$} & \multicolumn{2}{c}{Graded action cost} \\
\cmidrule(lr){2-3}\cmidrule(lr){4-5}
 & base & $+$critique & base & $+$critique \\
\midrule
0 & 0.430 & 0.851 & 544 & 268 \\
1 & 0.340 & 0.749 & 901 & 599 \\
2 & 0.471 & 0.871 & 535 & 261 \\
3 & 0.482 & 0.874 & 428 & 216 \\
4 & 0.459 & 0.868 & 537 & 274 \\
\midrule
Mean & 0.436 & \textbf{0.843} & 589 & \textbf{324} \\
\bottomrule
\end{tabular}
\end{table}

\begin{table}[!t]
\centering
\caption{System~2 duty cycle by evaluation protocol. The $1000$-window SisFall rate comes from the separate timing run, not the primary classification counts. SisFall protocols are stratified and D18/D19 enriched: these rates constrain feasibility but are neither free-living estimates nor guaranteed upper bounds. Distribution shift can change escalation in either direction. Binary KFall invokes deliberation on roughly one window in ninety under its own protocol.}
\label{tab:escalation}
\begin{tabular}{lcc}
\toprule
Protocol & $n$/fold & Escalation rate \\
\midrule
SisFall timing (5-fold mean)    & 1000 & 0.47 $\pm$ 0.08 \\
SisFall ambiguous stress        & 300  & 0.59 $\pm$ 0.04 \\
KFall external (5-fold mean)    & 1000 & 0.011 \\
\bottomrule
\end{tabular}
\end{table}

\begin{table}[!t]
\centering
\caption{Processing latency from the separate Ollama/Mistral agentic timing run (SisFall, $n{=}1000$/fold; GPU evaluation, not wearable hardware). The primary detector is evaluated for every window; the escalated path is dominated by local LLM decoding. Mean stack is the measured per-window average. Fold~2 has the slowest escalated path ($22.85$\,s), not the highest duty cycle (its rate is $0.460$, versus $0.546$ on Fold~4). Window acquisition, queueing and clinical response time are not established by these measurements.}
\label{tab:latency}
\begin{tabular}{crrr}
\toprule
Fold & Primary (ms) & Escalated (s) & Mean stack (s) \\
\midrule
0 & 8.96 & 12.81 & 5.68 \\
1 & 8.09 & 10.66 & 3.73 \\
2 & 8.52 & 22.85 & 10.52 \\
3 & 8.36 & 11.00 & 5.99 \\
4 & 8.89 & 10.86 & 5.93 \\
\midrule
Mean & 8.56 & 13.64 & 6.37 \\
\bottomrule
\end{tabular}
\end{table}


\end{document}
